\documentclass[a4paper,10pt]{article}
\newcommand\subdir{/home/koen/LaTeX-setup/}
\input{\subdir/LaTeX-files/imports.tex}
\input{\subdir/LaTeX-files/master-internship.tex}
\usepackage{\subdir/LaTeX-files/aastex}
\usepackage{natbib}
\bibliographystyle{\subdir/LaTeX-files/astroadstitle.bst}

%Set the title, Author and date
\title{Notes Week 21}
\author{Koen Essers}
\tutorial{Master Internship}
\date{\today}

%Set the header style
\header{\tutorialstring}

%Begin the document
\begin{document}
\setuplayout
\section{Dredge-up and Abundances}

To start, I want to plot the radius evolution for the models.
\begin{figure}[ht]
    \marginfignote{3}{All seem fine, and the general trends that I saw last week seem to hold for the lower metallicity star as well.}
    \centering
    \incplot{w21-radius}
\end{figure}

Now, more interestingly, let me plot the dredge-up efficiencies.

\begin{figure}[ht]
    \marginfignote{3}{It appears that dredge is slightly more efficient for lower metallicity stars. This is in line with literature predictions \citep{herwig2005, karakas1902}. Metallicity also seems to have an inverse relation with the number of thermal pulses that the star will experience.}
    \centering
    \incplot{w21-lambda}
\end{figure}

Below I plot the C/O-number ratio for the different models.

\begin{figure}[ht]
    \marginfignote{3}{Both the solar model and the \(Z=0.0883\) model look as expected, with the subsolar model achieving a slightly higher C/O-ratio, as expected.}
    \marginfignote{10}{The \(Z=0.00453\) model shows the strong bump that we saw in some of the earlier higher mass models as well.}
    \centering
    \incplot{w21-co}
\end{figure}

Since I have a lot of profiles for this star, we can try to work out what is happening\sidenote{Unfortunately though, I do not have all of the abundances for all of the profiles, so we need something else that shows the same behaviour.}
.

Let me start by plotting the Z-fraction at the surface of all profiles for the low-metallicity model.

\begin{figure}[ht]
    \marginfignote{0}{This looks a lot like the plot above, which is good, because that means we can use this \(Z\)-fraction as a way to understand the sudden jump.}
    \centering
    \incplot{w21-Z-surface}
\end{figure}


I really like \href{https://s-koen.github.io/figures/#/plots/master-internship/w21/dredge-up}{this interactive figure}. It shows:
\begin{enumd}
    \item  Repeated hydrogen shell burning followed by unstable helium burning.
    \item  The creation of metals during he-burning followed by dredge up.
    \item  The loss of the envelope
    \item The effect of the superadiabatic sub-surface layers on the metal abundances at the surface.
\end{enumd}

Clearly the abundance at the surface is not necessarily the abundance in most of the envelope. Below, I try to show that even the model with the sudden jump in surface C/O-ratio still has a well-behaved envelope abundance over time.

\begin{figure}[ht]
    \marginfignote{3}{The envelope abundance follows a much more predictable staircase-like pattern.}
    \marginfignote{8}{Since the superadiabatic subsurface layers contain very little mass, the envelope value should provide a much better estimate for the actual abundances that get transferred during RLOF and wind mass transfer.}
    \centering
    \incplot{w21-Z-envelope}
\end{figure}


\bibliography{master.bib}
\end{document}
